\section{Twelf}
\label{sec:twelf}
In previous sections we showed how to formalize
a programming language with a grammar for defining
its syntax, static semantic (type checking) rules that
defined which programs in a language are well-formed,
and dynamic semantics defining how to execute programs in the language.
A formal definition of a language enables the ability to
prove properties about the language.
We presented two important properties,
preservation and progress, that relate the static semantics
(compile-time analysis) with the dynamic semantics
(runtime behavior).

We showed how to prove properties of languages by
structural induction.
These proofs are typically are long, tedious,
and involve many cases.
As a result, it is easy to make a mistake while
writing such proofs.
Furthermore, the length and detail of these proofs means
it is easy for a human to miss mistakes while checking
the proof.
However, verifying and deriving proofs can (sometimes)
be done automatically.
Proof assistants/automated theorem provers such as Twelf~\cite{Twelf},
Coq~\cite{Coq}, and Isabelle~\cite{Isabelle} are software tools
that enable one to encode theory in their respective languages.
Depending on the power of these tools, they can
verify proofs of theorems or derive proofs
of theorems.
Using any of these tools requires a very large learning curve,
and explaining all of the details of how they work
is beyond the scope of this paper.
To sketch an idea of how they work,
the remainder of this section explains
the Twelf encoding of $\minilang$.
The entire Twelf encoding is available
online~\cite{twelf-tutorial}.
Further details on using Twelf can be found in other
Twelf tutorials~\cite{TwelfTutorial}.

\subsection{$\minilang$ Syntax in Twelf}
File \code{syntax.elf} from~\cite{twelf-tutorial} contains
the Twelf encoding of $\minilang$'s syntax.
In Twelf, there are three \emph{levels} of objects.
\emph{Kinds} are at the highest level.
\emph{Types} are at the second level.
\emph{Terms} are at the lowest level.
Each type is of a certain kind.
Each term is of a certain type.
For example, we could think of
the type \code{Array[Int]}, which is an
array of integers, to be of kind \code{Array}.
A term of type \code{Array[Int]} could be
\code{[1, 2, 3, 4]}.
In Twelf, one defines their language by defining
kinds, types, and terms.
The kind \code{type} is a primitive kind defined
in the Twelf language.
More information on Twelf's type system can be
found in~\cite{HARPER:2007:MML:1296837.1296842}.

Next, we describe the statements in file \code{syntax.elf}
using the terminology presented above.
On line 4,
``\code{exp :\ type}'' states that \code{exp} is a type
of kind \code{type}.
Similarly, line 7 defines \code{typ} to be type
of kind \code{type},
and line 10 defines \code{nat} to be a type
of kind \code{type}.
Line 11 defines the term \code{z} to be of
type \code{nat}.
Line 12 defines \code{s} to be a function term
of the function type from \code{nat}'s to \code{nat}'s.
Lines 11 and 12 encode (Peano) natural numbers
\code{\{z, s(z), s(s(z)), s(s(s(z))), \ldots \}}.
Lines 14--22 define strings to be a sequence of
the \code{char}'s separated by commas.
The functions \code{enat} and \code{estr}
are just ways to say that natural
numbers and strings are also expressions.
For example, \code{z} has type \code{nat},
but \code{(enat z)} has type \code{exp}.
\footnote{It would be nice if Twelf had the notion of
subtyping.
Then we could just tell Twelf \code{nat} $<:$ \code{exp},
so Twelf would infer that \code{z} should also have type
\code{exp} without having to wrap it as \code{(enat z)}.}

Lines 29--32 define the category of expressions
with function types.
For instance,
``\code{add :\ exp -> exp -> exp.}''
says that \code{add} is a function that takes into
two \code{exp}'s and returns an \code{exp}.
This corresponds to the grammar rule
\[\mathit{Expression} \; e \; \bnfdef \; \code{+($e_1$; $e_2$)}\]
stating that the addition of two expressions is also an expression.
\footnote{There are a few minor things missing in the definition
of $\minilang$ in this paper and its encoding in Twelf.
For example, line 32 defines an expression for representing the
length of a string (\code{len($e$)}).
Some details were left out for brevity.}
Similarily, ``\code{cat :\ exp -> exp -> exp.}''
says that \code{cat} is a function that takes into
two \code{exp}'s and returns an \code{exp};
\code{cat} corresponds to the string concatenation of two
subexpressions.
The last two lines of \code{syntax.elf} define two terms
\code{num} and \code{string} that represent the only two
types in $\minilang$.

The number of arguments a function term defined in Twelf
takes may not be clear to a reader who is not familiar with
\emph{currying}~\cite{CurryHaskel,CurryScala}.
For example, the \code{add} function term looks like a
function that takes in a single \code{exp} term and returns
another function of type \mbox{\code{exp -> exp}}.
Functions in Twelf are applied to terms by currying,
where there is no distinction between
functions $f$ and $f'$ when passing them two arguments
if $f(x)$ returns another function $g(y)$ and
\mbox{$f'(x, y)$ = $\underbrace{f(x)}_{g}(y)$}.
Passing a function multiple arguments is transformed
into a chain of function calls, where each function in
the chain is applied to a single argument.
So a function that returns another function
can be thought of as a function that takes in multiple
arguments.
Conversely, if a function $f$ takes in multiple
arguments $x_1, x_2, \ldots, x_n$ (where $n \geq 2$),
then $f(t)$ can be thought of as a function $g$
that is the same as $f$ except that occurrences of variable $x_1$
in the body of $f$ are replaced with term $t$ in the body of $g$.

\subsubsection{Representing terms w/ variables in Twelf
using Higher-Order Abstract Syntax}
\label{sec:twelf-hoas}

The encoding of the \texttt{let} term in Twelf uses the
technique of \emph{higher-order abstract syntax} (HOAS)~\cite{Twelf-HOAS}.
HOAS is a technique for representing abstract syntax trees
with bound variables.
The abstract syntax from Section~\ref{sec:grammar} used to describe
the grammar of $\minilang$ is actually
\emph{first-order abstract syntax} (FOAS).
In FOAS, each AST has the form $o(t_1, t_2, \ldots, t_n)$,
where $o$ is an operator and $t_1, t_2, \ldots, t_n$ are each
AST themselves.
The operands of an operator correspond to subexpressions.
For example, in ``\texttt{+($e_1$; $e_2$)}'',
$e_1$ and $e_2$ are operands/subexpressions of the
\texttt{+} operator.
The Twelf encoding of ``\texttt{+($e_1$; $e_2$)}''
is ``\texttt{add : exp -> exp -> exp.}'',
where the \texttt{add} corresponds to \texttt{+},
the leftmost \texttt{exp} corresponds to $e_1$,
and the middle \texttt{exp} corresponds to $e_2$.

In HOAS, ASTs declare the variables they bind.
For example, consider the AST
$o(t_1, t_2, \ldots, t_n)$.
In HOAS, each $t_i$ has the form
$x_1, x_2, \ldots x_k.t$, where $t$ is an AST,
each $x_j$ is a variable bound in $t$, and $k \geq 0$;
if $k = 0$, then $t_i$ does not introduce new variables in $t$.
One advantage of knowing each variable introduced by an AST
is that we can rename variables without changing the meaning
the AST.

Now consider the \texttt{let} expression in $\minilang$ and
its representation in FOAS:
\[\code{let($x$; $e_1$; $e_2$)}\]
Recall that this \code{let} expression is evaluated by
binding the variable $x$ to the value of $e_1$ in $e_2$.
Hence, the \code{let} expression can be represented by the
following HOAS:
\[\code{let($e_1$; $x.e_2$)}\]
The ``$x.e_2$'' captures that $x$ is bound in $e_2$.
HOAS lets us know where variables are being bound.
This lets us easily determine that the following two expressions
are equivalent, since they only differ in variable names:
\[ \code{let(3; $x$.+($x$; 4))} \equiv \code{let(3; $y$.+($y$; 4))}\]

Our Twelf encoding of the \code{let} expression encodes its
HOAS version, which can be seen with the type of its
second argument: \code{(exp -> exp)}.
Usually, each operator such as \code{add} and \code{cat}
only took in \code{exp} arguments, where the arguments
represented subexpressions.
However, the second argument to \code{let} is not just an
ordinary subexpression but a subexpression that introduces
a new variable.
The second argument to \code{let} represents a higher-order
AST of the form ``$x.e_2$''.

An AST that introduces new variables is represented in
Twelf by a function.
Functions in programming languages are really just terms
with \emph{holes}.
The holes are represented by free variables,
and these holes are filled in when these
(terms w/ holes)/functions are applied to other terms.
Hence, ``$x.e_2$'' can be thought of as the function
or \emph{lambda-abstraction} ``$\lambda x.e_2$'',
where $e_2$ is the body of the function and $x$ is a variable
that is bound to a term in $e_2$.
A term $t$ of type \code{(exp -> exp)} will be a term of type
\code{exp} that contains occurrences of a free variable that
will be replaced with another term when $t$ is applied to it.
Twelf's syntax for the function ``$\lambda x.e$'' is
``\code{[$x$] $e$}''.
Figure~\ref{fig:hoas} gives a concrete example of a
\texttt{let} expression in
concrete syntax and its corresponding representation
in Twelf's HOAS.

\begin{figure}[H]
\begin{center}
\begin{tabular}{l|l}
\textit{Concrete Syntax} & \textit{Twelf HOAS} \\ \hline
\code{let x = $\underbrace{\code{1 + 2}}_{e_1}$ in
                $\underbrace{\code{x + 3}}_{e_2}$} &
  \code{let $\underbrace{\code{(add 1 2)}}_{e_1}$
    $\underbrace{\code{([x] add x 3)}}_{x.e_2}$}
\end{tabular}
\end{center}
\caption{Example Higher-Order Abstract Syntax in Twelf}
\label{fig:hoas} 
\end{figure}




\subsection{\texorpdfstring{$\minilang$}{MiniLang} Static Semantics and Judgments in Twelf}
File \code{typing.elf} from~\cite{twelf-tutorial} contains the encoding of
$\minilang$'s static semantics in Twelf.
Line 5, ``\code{of :\ exp -> typ -> type.}'',
defines the relation \code{of} representing
the judgment $e : \tau$.
It defines \code{of} to be a function kind or
\emph{type family}~\cite{GHCTypeFamily,TwelfTypeFamily}.
Type families are functions that return types instead of terms.
The ``\code{of}'' relation can be thought of as (1) a function
that returns types when applied to \code{exp} and \code{typ} terms
or (2) as a set of types \emph{indexed} by \code{exp} and \code{typ} terms.

Judgments are represented in Twelf as
\emph{dependent types}~\cite{TwelfTypeJudge,advancetypes_ch02,Twelf-DependentTypes}.
The \code{of} type family returns dependent types of kind \code{type}.
Hence, function \code{of} returns judgments.
% Introduce dependent types
Dependent types are types that include terms as one of its components.
Example dependent types come from the set of $n$-dimensional vectors
of real numbers denoted \code{Vec($n$)}.
For instance, suppose \code{3} is a term representing the number $3$,
$\nats$ is a type representing the set of natural numbers,
and \code{3} is of type $\nats$.
Also, assume \code{Vec} is a function that takes in a natural
number $n$ as input and returns a type representing
the set of $n$-dimensional vectors.
Then, \code{Vec(3)} is a dependent type representing
$3$-dimensional vectors of real numbers.
A term of type \code{Vec(3)} is the vector $[7, 3, 4]$.
Furthermore, function \code{Vec} denotes a type family
indexed by natural numbers:
\code{Vec(0)}, \code{Vec(1)}, \code{Vec(2)}, \code{Vec(3)}, \ldots

% Back to Twelf
The \code{of} type family is indexed by \code{exp} and \code{typ} terms.
The dependent type ``\mbox{\code{of $e$ $\tau$}}'' represents
the judgment $e : \tau$.
For example, dependent type ``\mbox{\code{of (enat z) num}}''
represents the proposition or judgment $z :\ \code{num}$.
A proof or derivation of a judgment/type is represented by a term of that type.
An example proof will be presented later in this section.

Inference rules in Twelf are modeled as functions that return terms
of dependent types.
Consider the function term \code{of/nat} defined on line 7
of file \code{typing.elf}.
Strings that begin with capital letters in Twelf are interpreted
to be parameters.
The \code{N} parameter in ``\mbox{\code{of (enat N) num}}'' is passed
to the \code{enat} function term.
Since \code{enat} only takes in values of type
\code{nat}, so does \code{of/nat}.
Twelf is able to infer that the \code{N} variable must be
bound to a term of type \code{nat}.

At first, it seems the \code{of/nat} function returns types.
For example, it seems ``\code{of/nat z}'' returns
the dependent type ``\mbox{\code{of (enat z) num}}''.
However, ``\code{of/nat z}'' actually returns
a \emph{term} of type ``\mbox{\code{of (enat z) num}}''.
The term returned by ``\code{of/nat z}'' represents a derivation
or proof of the judgment $z :\ \code{num}$;
this judgment is represented by type ``\mbox{\code{of (enat z) num}}'',
which is the type of term ``\code{of/nat z}''.
Hence, the \code{of/nat} term corresponds to rule \code{T.1},
which is the rule that allows us to derive that any natural
number has type \code{num}.
Similarly, the \code{of/str} function term takes in any string $s$
and returns a derivation/term of a judgment/type stating that
string $s$ has type \code{string}.

Function term \code{of/add} models rule \code{T.4}.
Premises of inference rules are represented by inputs that must be
terms of a dependent type.
For example, function \code{of/add} takes in two
(explicit~\cite{Twelf-Implicit}) arguments of dependent types.\footnote{
Function \code{of/add} actually takes in four arguments.
The first two arguments are the \code{exp} terms \code{E1} and
\code{E2}.
However, these arguments are \emph{implicit} arguments and
typically do not need to be mentioned because they can be inferred
from the last two \emph{explicit} arguments of dependent types.
For example, since the dependent type ``\mbox{\code{of E1 num}}'' includes
\code{E1} as one if its components, Twelf can extract \code{E1}
from this type.
For more information on implicit and explicit arguments
see~\cite{Twelf-Implicit}.}
Variables \code{E1} and \code{E2} are bound to terms of type \code{exp}.
Dependent type ``\mbox{\code{of E1 num}}'' represents the judgment
$e_1 : \code{num}$.
A term of type ``\mbox{\code{of E1 num}}'' represents a derivation
of $e_1 : \code{num}$.
Hence, \code{of/add} takes in a derivation of $e_1 : \code{num}$
and a derivation of $e_2 : \code{num}$ and returns
a derivation/term of type ``\mbox{\code{of (add E1 E2) num}}'';
this type represents judgment \mbox{\code{+($e_1$; $e_2$) $:$ num}}.

\subsubsection{Higher-Order Judgments in Twelf}
Inference rules involving typing contexts
such as rule \code{T.6} are \emph{hypothetical judgments}~\cite{TwelfHypothetical}
or judgments that make use of hypothetical assumptions.
Hypothetical judgments can be encoded in Twelf using
\emph{higher-order judgments}~\cite{TwelfHigherOrder}
or judgments that contain other judgments.
Higher-order judgments are encoded in Twelf as \emph{higher-order types}.
Higher-order types are functions types, where one of their input types
is also a function type.
Encoding higher-order judgments as higher-order types is similar to
the technique of higher-order abstract syntax discussed in
Section~\ref{sec:twelf-hoas}, where syntactic terms with binders
(e.g. the body of the \code{let} expression) are encoded as
terms with holes or equivalently
functions that take in terms as input and return terms.
Higher-order judgments are encoded with function types that have
input types representing the hypothetical assumptions.

Higher-order types representing hypothetical judgments involve
\emph{pi-types}, denoted $\Pi x:S.T$.
A pi-type is a special kind of function type.
Terms of pi-types are \emph{pi-abstractions}.
First, a \emph{lambda-abstraction}, denoted ``$\lambda x:S.e$'',
is a function term of a more conventional kind of function type,
 $\ftype{S}{T}$.
A function of this type takes in a term of type $S$ and
returns a term of type $T$.
However, a pi-abstraction of pi-type $\Pi x:S.T$ would map
a term $s$ of type $S$ to a term of type $[s / x]T$.
That is, the return type of a pi-type can vary according
to the argument supplied.
Hence, if $x$ is (syntactically) a part of $T$ in the pi-type
$\Pi x:S.T$, then a term returned by a function of that
pi-type is a term of a dependent type, since the return type
$T$ depends on the input term $x$.
If $x$ is not a part of $T$ in $\Pi x:S.T$,
then we abbreviate $\Pi x:S.T$ as $\ftype{S}{T}$,
which also signals that a lambda-abstraction can be of
this type, since the return type does not involve the argument.
Lastly, the pi-type $\Pi x:S.T$ is represented in Twelf syntax
as \code{\{x:S\} T}.

A common Twelf coding convention is used in
the return type of the \code{of/let} function term
given on line 20 of the \code{typing.elf} file:
``\mbox{\code{of (let E1 ([x] E2 x)) T2}}''.
This code snippet could have been replaced with
the shorter snippet:
``\mbox{\code{of (let E1 E2) T2}}''.
The \code{E2} in the shorter snippet
already denotes a term
of function type ``\mbox{\code{exp -> exp}}''
(or equivalently, an \code{exp} with free variable
occurrences).
However, we applied a Twelf coding convention when
we replaced \code{E2} with the equivalent function
term \code{([x] E2 x)}.
The wrapper function \code{([x] E2 x)} is used 
just to make it easier to see that \code{E2} is
a function term that takes in a single argument.

The \code{of/let} term represents rule \code{T.6}
and is of a higher-order type:
\begin{verbatim}
of/let :  ({x: exp} of x T1 -> of (E2 x) T2) ->
          of E1 T1 ->
          of (let E1 ([x] E2 x)) T2
\end{verbatim}
The second argument type ``\mbox{\code{of E1 T1}}''
corresponds to the premise ``$e_1 : \tau_1$''.
The type of \code{of/let}'s first argument is the pi-type:
``\mbox{\code{\{x: exp\} of x T1 -> of (E2 x) T2}}'';
let $P$ denote this pi-type.
A pi-abstraction of type $P$ takes in two arguments:
The first is an \code{exp} term, which will be bound to
\code{x}.
The second is a term of type ``\mbox{\code{of x T1}}''
representing a proof that \code{x} has type \code{T1}.
Given two such terms,
a pi-abstraction of type $P$ should return
a derivation of ``\mbox{\code{of (E2 x) T2}}'',
where \code{(E2 x)} is the \code{E2} term with its hole
filled in by the \code{exp} term that is bound to \code{x}.

The pi-type $P$ represents the second premise
``\mbox{$\Gamma , x : \tau_1 \vdash e_2 : \tau_2$}''
of rule ~\code{T.6}.
Passing a derivation \code{dx} of type ``\mbox{\code{of x T1}}''
to a pi-abstraction simulates extending the typing context
with the assumption ``\mbox{$x : \tau_1$}''.
The derivation \code{dx} of type ``\mbox{\code{of x T1}}''
can be used in the body of the pi-abstraction to
return a term/proof of ``\mbox{\code{of (E2 x) T2}}''.


\subsection{Twelf Theorems, Proofs, and Wrap-up}
Explaining all of the details of the Twelf code is
beyond the scope of this paper.
So in this section, we conclude the Twelf discussion with a brief
overview of the remaining files of the Twelf $\minilang$ encoding,
how theorems are specified in Twelf, example Twelf proofs,
and how Twelf helps with proving theory about languages.


\subsubsection{Twelf Project Files}
Twelf projects are stored in directories containing a configuration
file typically named \code{sources.cfg}.
The \code{sources.cfg} file tells Twelf the files to read
and the order in which to process them.
The \code{evaluation.elf} file contains the dynamic semantics
of $\minilang$.
File \code{preservation.elf} contains the preservation theorem
and its proof.
File \code{progress.elf} contains the progress theorem
and its proof.
The other files are explained throughout Section~\ref{sec:twelf}.

\subsubsection{Theorems are Function Types. Proofs of Theorems are Function Terms.}
Theorems in Twelf are encoded similarly to
(for all)-(there exists) queries.
For example, below is the encoding of the progress theorem
from lines 11--12 and 153--154 of \code{progress.elf}
(shown without an inline comment).

\begin{verbatim}
progress : of E T -> not_stuck E -> type.
%mode progress +O -NS.
...
%worlds () (progress O NS).
%total O (progress O _).
\end{verbatim}

The first line declares \code{progress} as a type family:
a relation between a derivation of the typing judgment
\code{of E T} and a derivation of the judgment that expression
\code{E} is not stuck (\code{not\_stuck E}).
The \code{\%mode} declaration marks the derivation \code{O}
of \code{of E T} as an input (\code{+}) and the derivation
\code{NS} of \code{not\_stuck E} as an output (\code{-});
it states that for every derivation \code{O}, there exists
a derivation \code{NS}.
The clauses between these declarations, each beginning with
``\code{- :}'', are the cases of the proof.
The \code{\%worlds} declaration states the contexts in which
the theorem holds; the empty parentheses mean that it is only
stated for closed terms, with no hypothetical assumptions.
Finally, \code{\%total} asks Twelf to check that the relation
is \emph{total}:
that the cases cover every possible input derivation
and that the proof terminates, by induction on \code{O}.
If the check succeeds, Twelf has verified the proof.
Twelf also has a more verbose syntax for stating theorems,
the \code{\%theorem} declaration; \code{preservation.elf}
shows the preservation theorem in that syntax in a comment.
For progress, it would be written as follows.

\begin{verbatim}
%theorem
progress :
  forall* {E} {T}
  forall {O:of E T} exists {NS:not_stuck E}
  true.
\end{verbatim}

The \code{forall*} line declares the variables in the theorem.
The next line states that for every derivation of the typing judgment,
\code{of E T}, where the derivation of that judgment is bound to \code{O}, there
exists a derivation of the judgment that expression \code{E} is
not stuck (\code{NS:not\_stuck E}).
The following two rules define the \code{not\_stuck} judgment.\footnote{
In Twelf, function types can be written using arrows
pointing either to the left or to the right.}

\code{not\_stuck/val: not\_stuck E <- value E.} states that if
\code{E} is a value, then \code{E} is not stuck.

\code{not\_stuck/step: not\_stuck E <- step E E'.} states that if
there exists an \code{E'} such that \code{step E E'},
then \code{E} is not stuck as well.

Note that the above progress theorem is nothing more than a function
type.
Given a term/proof of the judgment \code{of E T},
a function of this type should return a term/proof of
the judgment \code{not\_stuck E}.
Hence, a proof of the progress theorem/(function type) is a (total)
function or term of the function type representing the theorem.

\subsubsection{Checking and Deriving Proofs}
The benefit of using Twelf is that it can check your proofs
and occasionally derive proofs.
File \code{test\_typing.elf}, for example, provides two
example judgments that can be derived automatically by Twelf.

\begin{verbatim}
%solve D1 : of (estr (a , b , c , a , eps)) string.
%solve D2 : of (add (enat (s z)) (enat z)) num.
\end{verbatim}

A \code{\%solve} declaration asks Twelf to search for a term
of the given type and, if it finds one, to define a new
constant with that term as its value.
The first declaration asks Twelf to derive that the type of
\code{str[abca]} is \code{str},
and save that derivation in \code{D1}.
The second declaration asks Twelf to derive that the
type of \code{+(num[1]; num[0])} is \code{num},
and save that derivation in \code{D2}.

Below is the Twelf output from those declarations.

\begin{verbatim}
[Opening file test_typing.elf]
%solve
of (estr (a , b , c , a , eps)) string.
 OK
D1 : of (estr (a , b , c , a , eps)) string = of/str.
%solve
of (add (enat (s z)) (enat z)) num.
 OK
D2 : of (add (enat (s z)) (enat z)) num = of/add of/nat of/nat.
[Closing file test_typing.elf]
\end{verbatim}

Twelf finds derivations of both judgments.
It prints them with implicit arguments omitted
(see the footnote in Section~\ref{sec:twelf-static}).
Written with all of their arguments, the derivations are
\begin{verbatim}
D1 = of/str (a , b , c , a , eps).
D2 = of/add (enat z) (enat (s z)) (of/nat z) (of/nat (s z)).
\end{verbatim}
Derivation \code{D1} says that applying function \code{of/str}
to the string term
\code{(a , b , c , a , eps)} returns a
term/derivation of type/judgment
``\code{of (estr (a , b , c , a , eps)) string}''.

Derivation \code{D2} is described as follows:
\begin{enumerate}
\item Apply function \code{of/nat} to term \code{z}
to return a term/derivation of ``\code{of (enat z) num}''.
\item Apply function \code{of/nat} to term ``\code{s z}''
to return a term/derivation of ``\code{of (enat (s z)) num}''.
\item Apply \code{of/add} to the derivations
of ``\code{of (enat z) num}'' and ``\code{of (enat (s z)) num}''
to return a derivation of
``\code{of (add (enat (s z)) (enat z)) num}''.
The first argument \code{(enat z)} establishes
that the third argument should be a term of type
``\code{of (enat z) num}'';
hence, the type of the third argument of \code{of/add}
depends on the first argument passed to \code{of/add}.
The situation is analogous for the second and fourth
arguments of function \code{of/add}.
\end{enumerate}

Twelf could not derive the proof of the preservation theorem,
but it could let you know if your proof was incorrect.
For example, lines 23--27 in \code{preservation.elf}
prove preservation for the typing-evaluation rule combination
(T.4, D.1).
If those lines were removed, Twelf would return the following
error output:

\begin{verbatim}
preservation.elf:98.8-98.11 Error: 
Coverage error --- missing cases:
{E1:exp} {E2:exp} {E3:exp} {O1:of (add E1 E2) num} {S1:step E1 E3}
   {O2:of (add E3 E2) num}
   |- preservation O1 (step/add1 S1) O2.
\end{verbatim}

The error message lets us know that we forgot to prove preservation
for the case when we could derive ``\code{of (add E1 E2) num}''
and ``\code{step (add E1 E2) (add E3 E2)}''
(\code{+($e_1$; $e_2$)} $\stepto$ \code{+($e_3$; $e_2$)}).
In this case, judgment ``\code{step (add E1 E2) (add E3 E2)}''
was derived by applying function (inference rule) \code{step/add1}
to a derivation (\code{S1}) of type ``\code{step E1 E3}''.
To prove preservation for this case,
we would need to produce a derivation of type ``\code{of (add E3 E2) num}''.
This would show that performing an evaluation step on the left
summand of a well-typed \code{add} expression would not change
the type of the resulting \code{add} expression.

In summary, Twelf can aid in deriving simple judgments and checking
that proofs of language properties are correct.
If a proof is not correct, Twelf will let us know that the proof
contains a mistake and provide an error message to help ``debug'' the proof.


